\begin{proposition}[A surface-support criterion for the first Chow gate]
\label{prop:B-surface-supported-CH2}
\status{Proved; the unrestricted threefold gate remains unresolved}
Let \(X/k\) be smooth, projective and geometrically integral, where
\(k\) is finite, and put \(n=\dim X\geq2\).
Suppose that, after a finite extension of \(k\), there is a closed
subscheme \(S\subset X\), of dimension at most two, for which
\[
 [\eta_X]\in\operatorname{im}\bigl(
 CH_0(S_{k(X)})_{\Q}\longrightarrow
 CH_0(X_{k(X)})_{\Q}\bigr).
 \tag{B.S1}\label{eq:B-generic-surface-support}
\]
Then rational and numerical equivalence agree on codimension-two
cycles on \(X\).  In particular, for every \(\ell\ne\operatorname{char}k\),
\[
 CH^2(X)_{\Q}\otimes_{\Q}\Q_\ell\longrightarrow
 H^4(X_{\overline k},\Q_\ell(2))
 \quad\hbox{is injective}.
\]
Condition~\eqref{eq:B-generic-surface-support} is equivalent to
universal rational support on this fixed \(S\): the pushforward
\(CH_0(S_L)_{\Q}\to CH_0(X_L)_{\Q}\) is surjective for every
field extension \(L/k\).  Thus it is a condition over function
fields, not merely a condition on zero-cycles over the finite base
field.  For threefolds satisfying it, rational and numerical
equivalence agree in every ordinary Chow degree.
\end{proposition}
\begin{proof}
All assertions descend through a finite constant extension by
restriction and corestriction, so make the extension in the hypothesis.
Clearing denominators in~\eqref{eq:B-generic-surface-support},
spreading its finitely many cycles and rational functions over a
nonempty open of the first factor, and using Chow localization gives
\[
 N\Delta_X=\Theta_1+\Theta_2\quad\hbox{in }CH^n(X\times X),
 \qquad
 \operatorname{supp}\Theta_1\subset X\times S,
 \quad\operatorname{supp}\Theta_2\subset D\times X,
\]
where \(N>0\) and \(D\subset X\) is a divisor containing the
closed complement of that open.  This is the generic-point proof
of the Bloch--Srinivas decomposition, so the support implication
does not require an additional conjecture.  The corresponding
curve-support factorization is explicitly proved in
\citet[Lemma~5.1 and its proof, p.~841]{BFuVial2023}.
Conversely, after any extension \(L/k\), move a zero-cycle on the
smooth variety \(X_L\) away from \(D_L\).  The action of
\(\Theta_2\) is zero on that cycle, while the action of
\(\Theta_1\) lands in the image of \(CH_0(S_L)\).
Division by \(N\) proves universal rational support.  The other
implication follows by taking \(L=k(X)\) and the generic point.

Transpose the decomposition and write
\[
 N\Delta_X=\Gamma_1+\Gamma_2,
 \qquad
 \operatorname{supp}\Gamma_1\subset S\times X,
 \quad\operatorname{supp}\Gamma_2\subset X\times D.
 \tag{B.S2}\label{eq:B-transposed-surface-support}
\]
Replace the irreducible components of \(S\) and \(D\) by smooth
projective alterations, separately when necessary
\citep[Theorem~4.1 and Remark~4.2]{c_dejong1996}.  Denote their
maps to \(X\) by \(a:S'\to X\) and \(b:D'\to X\), with
\(\dim S'=s\leq2\) and \(\dim D'=n-1\) on each component.
Proper surjective pushforward on Chow groups is surjective with
rational coefficients: for an integral cycle downstairs, choose a
closed point of its generic fibre, take its closure upstairs, and
divide by the resulting finite generic degree.  Consequently the
two supported correspondences lift, component by component, to
\[
 \widetilde\Gamma_1\in CH^{s}(S'\times X)_{\Q},
 \qquad
 \widetilde\Gamma_2\in CH^{n-1}(X\times D')_{\Q}.
\]
Their actions on codimension-two Chow groups have the precise
factorizations
\[
\begin{aligned}
 (\Gamma_1)_*&=(\widetilde\Gamma_1)_*a^*:
 CH^2(X)_{\Q}\longrightarrow CH^2(S')_{\Q}
                 \longrightarrow CH^2(X)_{\Q},\\
 (\Gamma_2)_*&=b_*(\widetilde\Gamma_2)_*:
 CH^2(X)_{\Q}\longrightarrow CH^1(D')_{\Q}
                 \longrightarrow CH^2(X)_{\Q}.
\end{aligned}
\]
The pullback \(a^*\) exists as the refined pullback for a
morphism between smooth varieties; the dimension shifts above
also follow directly by composing the graph correspondences.
All these maps preserve numerical equivalence, by the projection
formula, and preserve homological equivalence.

Here are the elementary finite-field vanishing inputs.  On any
smooth projective geometrically integral variety \(Y\) over a
finite field, every degree-zero zero-cycle is torsion: choose a
smooth projective geometrically integral curve through its finite
support, using \citet[Corollary~3.4, with Theorem~3.3 and its
proof]{DPoonen2004Bertini}; its class on that curve belongs to the
finite group of rational points of the Jacobian.  Pushing forward
annihilates a positive multiple of the original cycle.  A finite
constant extension handles geometric components if needed.
Thus the numerically trivial part of \(CH^2(S')_{\Q}\) is zero:
it is zero by dimension when \(s<2\), and is the degree-zero
part of \(CH_0(S')_{\Q}\) when \(s=2\).  Numerically trivial
divisors on \(D'\) are algebraically trivial up to torsion; the
remaining Picard-zero points over a finite field are torsion too.
Hence \(CH^1(D')_{\Q,\mathrm{num}}=0\).  The divisor assertion
is also the exact input in the proof of Fu--Vial's Lemma~5.1;
here no hypothesis on \(H^3(X)\) is needed.

For numerically trivial \(z\in CH^2(X)_{\Q}\), both actions vanish.
Equation~\eqref{eq:B-transposed-surface-support}
gives \(Nz=0\), hence \(z=0\).  Homologically trivial cycles
are numerically trivial.  To obtain the stated injectivity even
after tensoring with \(\Q_\ell\), recall that numerical Chow
groups are finite dimensional: a full-rank matrix of rational
intersection numbers embeds its rank in the finite-dimensional
cohomological pairing.  Choose a rational basis for \(CH^2(X)_{\Q}\)
and complementary cycles giving an invertible rational intersection
matrix.  That same matrix remains invertible over \(\Q_\ell\),
so the realized basis remains linearly independent over \(\Q_\ell\).
For a threefold the remaining ordinary degrees are degree zero,
divisors, and zero-cycles, already covered by the preceding
elementary statements.
\end{proof}

\begin{corollary}[Geometrically uniruled threefolds]
\label{cor:B-uniruled-threefold}
\status{Proved ordinary Chow equality; no full motivic comparison inferred}
Every smooth projective geometrically uniruled threefold over a
finite field satisfies rational equals numerical equivalence in
every ordinary Chow degree.  In particular the first unresolved
codimension-two threefold test in
Proposition~\ref{prop:B-diagonal-Chow-gate} may be restricted to
threefolds which are not geometrically uniruled.
\end{corollary}
\begin{proof}
After a finite constant extension, geometric uniruledness gives a
dominant rational map \(\phi:T\times\mathbf P^1\dashrightarrow X\),
where \(T\) is an integral surface.  The two sides have the same
dimension, so the map is generically finite.  After another finite
constant extension choose \(a\in\mathbf P^1(k)\) for which
\(\phi\) is defined at the generic point of \(T\times\{a\}\).
Let \(S\) be the closure of its image; then \(\dim S\leq2\).

Over \(K=k(X)\), choose a closed point in the generic fibre of
the domain of \(\phi\).  Its residue field \(K'/K\) is finite,
and the corresponding point of \(T\) is generic: a proper closed
subset of \(T\), multiplied by \(\mathbf P^1\), has dimension
at most two and cannot dominate \(X\).  The resulting rational
map \(\mathbf P^1_{K'}\dashrightarrow X_{K'}\) extends to a
morphism because the source is a regular curve and the target is
proper.  On this curve the selected \(K'\)-point and \(a\) are
rationally equivalent.  Their images show that
\([\eta_X]_{K'}\) is represented by a zero-cycle on \(S_{K'}\).
Corestriction proves~\eqref{eq:B-generic-surface-support} over
\(K\), after division by \([K':K]\).  The proposition applies,
and finite-constant descent finishes the proof.
\end{proof}

The support hypothesis must not be weakened to
\(CH_0(X_{\overline{\F}_p})_{\Q}=\Q\): that statement already
holds for every smooth projective geometrically integral
finite-field variety by the curve argument above.  Nor does the
proposition establish Parshin's higher-Chow vanishing, comparison
for all twists of an individual such variety, or the higher odd
dimensional tests in Proposition~\ref{prop:B-diagonal-Chow-gate}.
The criterion is not subsumed by our abelian-type theorem: for an
arbitrary smooth projective surface \(S\), \(S\times\mathbf P^m\)
has universal support on \(S\times\{a\}\), while its surface
Chow summand is not in general known to be of abelian type.
